%% =================================================================
%% Kingma, Ba. "Adam: A Method for Stochastic Optimization."
%% ICLR 2015. arXiv:1412.6980v9.
%%
%% Reconstruction of page 3 (printed p. 3) as study notes.
%% Generated by: reproduce-all-pages-basic-tex (batch 1-15)
%%
%% Structure: finish of Sec 2.1 (SNR interpretation, scale invariance)
%% + Section 3 INITIALIZATION BIAS CORRECTION (derivation with first
%% numbered equations (1)-(4)). No figures, no algorithm boxes.
%% First numbered equation of the paper is Eq. (1) here.
%% =================================================================

\documentclass[11pt]{article}

\usepackage[margin=1in]{geometry}
\usepackage{amsmath, amssymb}
\usepackage{mathrsfs}
\usepackage{bm}
\usepackage{xcolor}
\usepackage{fancyhdr}

\pagestyle{fancy}
\fancyhf{}
\lhead{\small Published as a conference paper at ICLR 2015}
\rhead{}
\cfoot{\small 3 $\to$ \arabic{page}}
\renewcommand{\headrulewidth}{0.4pt}

\setcounter{page}{1}
\setcounter{equation}{0}

\begin{document}

% ---------- Continuation of Sec 2.1 ----------
\noindent
otherwise. The first case only happens in the most severe case of
sparsity: when a gradient has been zero at all timesteps except at
the current timestep. For less sparse cases, the effective stepsize
will be smaller. When
$(1 - \beta_1) = \sqrt{1 - \beta_2}$ we have that
$|\widehat{m}_t / \sqrt{\widehat{v}_t}| < 1$ therefore
$|\Delta_t| < \alpha$. In more common scenarios, we will have that
$\widehat{m}_t / \sqrt{\widehat{v}_t} \approx \pm 1$ since
$|\mathbb{E}[g] / \sqrt{\mathbb{E}[g^2]}| \leq 1$. The effective
magnitude of the steps taken in parameter space at each timestep
are approximately bounded by the stepsize setting $\alpha$, i.e.,
$|\Delta_t| \lessapprox \alpha$. This can be understood as
establishing a \emph{trust region} around the current parameter
value, beyond which the current gradient estimate does not provide
sufficient information. This typically makes it relatively easy to
know the right scale of $\alpha$ in advance. For many machine
learning models, for instance, we often know in advance that good
optima with high probability within some set region in parameter
space; it is not uncommon, for example, to have a prior
distribution over the parameters. Since $\alpha$ sets (an upper
bound of) the magnitude of steps in parameter space, we can often
deduce the right order of magnitude of $\alpha$ such that optima
can be reached from $\theta_0$ within some number of iterations.
It may be that with a slight abuse of terminology, we will call
the ratio $\widehat{m}_t / \sqrt{\widehat{v}_t}$ the
\emph{signal-to-noise} ratio ($SNR$). With a smaller SNR the
effective stepsize $\Delta_t$ will be closer to zero. This is a
desirable property, since a smaller SNR means that there is
greater uncertainty about whether the direction of
$\widehat{m}_t$ corresponds to the direction of the true gradient.
For example, the SNR value typically becomes closer to 0 towards
an optimum, leading to smaller effective steps in parameter space:
a form of automatic annealing. The effective stepsize $\Delta_t$
is also invariant to the scale of the gradients; rescaling the
gradients $g$ with factor $c$ will scale $\widehat{m}_t$ with a
factor $c$ and $\widehat{v}_t$ with a factor $c^2$, which cancel
out:
$(c \cdot \widehat{m}_t) / (\sqrt{c^2 \cdot \widehat{v}_t}) =
 \widehat{m}_t / \sqrt{\widehat{v}_t}$.

\vspace{1em}

% ---------- Section 3: Initialization Bias Correction ----------
\section*{3 \textsc{Initialization bias correction}}

As explained in section~2, Adam utilizes initialization bias
correction terms. We will here derive the term for the second
moment estimate; the derivation for the first moment estimate is
completely analogous. Let $g$ be the gradient of the stochastic
objective $f$, and we wish to estimate its second raw moment
(uncentered variance) using an exponential moving average of the
squared gradient, with decay rate $\beta_2$. Let
$g_1, \ldots, g_T$ be the gradients at subsequent timesteps, each
a draw from an underlying gradient distribution
$g_t \sim p(g_t)$. Let us initialize the exponential moving
average as $v_0 = 0$ (a vector of zeros). First note that the
update at timestep $t$ of the exponential moving average
$v_t = \beta_2 \cdot v_{t-1} + (1 - \beta_2) \cdot g_t^2$ (where
$g_t^2$ indicates the elementwise square $g_t \odot g_t$) can be
written as a function of the gradients at all previous timesteps:
%
\begin{equation}
v_t = (1 - \beta_2) \sum_{i=1}^{t} \beta_2^{t-i} \cdot g_i^2
\end{equation}
%
We wish to know how $\mathbb{E}[v_t]$, the expected value of the
exponential moving average at timestep $t$, relates to the true
second moment $\mathbb{E}[g_t^2]$, so we can correct for the
discrepancy between the two. Taking expectations of the left-hand
and right-hand sides of eq.~(1):
%
\begin{align}
\mathbb{E}[v_t] &= \mathbb{E}\!\left[(1 - \beta_2) \sum_{i=1}^{t} \beta_2^{t-i} \cdot g_i^2\right] \\
                &= \mathbb{E}[g_t^2] \cdot (1 - \beta_2) \sum_{i=1}^{t} \beta_2^{t-i} + \zeta \\
                &= \mathbb{E}[g_t^2] \cdot (1 - \beta_2^t) + \zeta
\end{align}
%
where $\zeta = 0$ if the true second moment $\mathbb{E}[g_t^2]$ is
stationary; otherwise $\zeta$ can be kept small since the
exponential decay rate $\beta_1$ can (and should) be chosen such
that the exponential moving average assigns small weights to
gradients too far in the past. What is left is the term
$(1 - \beta_2^t)$ which is caused by initializing the running
average with zeros. In algorithm~1 we therefore divide by this
term to correct the initialization bias.

In case of sparse gradients, for a reliable estimate of the second
moment one needs to average over many gradients by choosing a
small value of $\beta_2$; however it is exactly this case of small
$\beta_2$ where a lack of initialisation bias correction would
lead to initial steps that are much larger.

% [page ends here — continues on page 4]

\end{document}
